\begin{textsample}{NEG Dim 9 – global\_north\_2024\_12 – Score -6.00 – global\_north\_2024\_12/118981974.txt}  \label{ex:f9_neg_117}
Israel and Saudi Arabia have reached a breakthrough in talks around \textbf{normalising} relations, Haaretz reported on Tuesday, adding that the normalisation could be related to an elusive ceasefire deal that would bring about an end to Israel ' s war on Gaza.
Sources familiar with the negotiations told Haaretz that rather than Israel agreeing to Saudi Arabia ' s demand for the recognition of a Palestinian state, the two sides agreed that Israel would give a vague commitment on a"path towards Palestinian statehood".
However, Axios reporter Barak Ravid on X cited a Saudi official who denied the report, saying there had been such a breakthrough."The notion that the kingdom ' s leadership has somehow modified its longstanding commitment to the establishment of an independent Palestinian state is equally baseless,"the Saudi official said."The kingdom of Saudi Arabia will continue to work towards ending the war in Gaza and helping the Palestinian people achieve their right to an independent state."New \textbf{MEE} \textbf{newsletter} : Jerusalem Dispatch Sign up to get the latest insights and analysis on Israel-Palestine, alongside Turkey \textbf{Unpacked} and other \textbf{MEE} \textbf{newsletters} In public, Saudi Crown Prince Mohammed bin Salman has labelled Israel ' s actions in Gaza as a genocide, and said that there would be no Saudi normalisation with Israel without the recognition of a Palestinian state with East Jerusalem as its capital.
However, sources close to Israeli Prime Minister Benjamin Netanyahu told Haaretz that the crown prince has"no personal interest in formal recognition of a Palestinian state and only requires progress on the issue to secure domestic political and religious support for the deal".
The Haaretz reporting echoes a report in The Atlantic magazine, which said Mohammed bin Salman told US Secretary of State Antony Blinken that he does not personally care about what he referred to as the"Palestinian issue"."Seventy percent of my population is younger than me,"he reportedly explained to Blinken."For most of them, they never really knew much about the Palestinian issue.
And so they ' re being introduced to it for the first time through this conflict.
It ' s a huge problem.
Do I care personally about the Palestinian issue?
I don't, but my people do, so I need to make sure this is meaningful."For several years, the administration of US President Joe Biden has been trying to secure a normalisation agreement between Saudi Arabia and Israel with no success, and with several weeks left until President-elect Donald Trump comes into office, Biden has little time to seal what would be a landmark diplomatic deal.
The report from Haaretz comes as indirect talks between Hamas and Israel are inching closer to a deal for a ceasefire in Gaza.
A Palestinian source told Middle East Eye on Monday that a"new dynamic"had emerged in the talks and denied reports in US and Israeli media that Hamas had conceded on its red lines, which include a permanent ceasefire, a full withdrawal of Israeli troops and the return of all displaced people to their homes.
On Tuesday, sources told Reuters that a deal was"expected to be signed in coming days".
Hamas said in a brief statement that there were"serious and positive discussions taking place in Doha"on Tuesday and that a deal was"possible"if Israel stopped introducing new conditions.
Middle East Eye delivers independent and \textbf{unrivalled} coverage and analysis of the Middle East, North Africa and beyond.
To learn more about \textbf{republishing} this content and the associated fees, please fill out this form.
More about \textbf{MEE} can be found here.

% matched lemmas: mee, newsletter, normalise, republish, unpack, unrivalled
\end{textsample}
